\chapter{Series and Parallel Networks}
\label{ch:resistive}
Resistive networks set bias points, divide voltages, sample currents, terminate
transmission lines, and model losses. They are the simplest place to practice the
conservation laws of \cref{ch:foundations}---KCL and KVL---and the
equivalent-circuit idea, so every rule below is derived from those laws rather than
memorized. Because the same two laws also fix how capacitors and inductors combine,
this chapter collects all three components' series and parallel rules in one place.

\section{Series and Parallel from KCL and KVL}
\textbf{Series} resistors carry the \emph{same current} \(I\). KVL sums their
voltages, so
\[
V=IR_1+IR_2+\cdots=I\sum_k R_k
\quad\Longrightarrow\quad
R_{\mathrm{eq}}=\sum_k R_k .
\]
\textbf{Parallel} resistors share the \emph{same voltage} \(V\). KCL sums their
currents, so
\[
I=\frac{V}{R_1}+\frac{V}{R_2}+\cdots=V\sum_k\frac{1}{R_k}
\quad\Longrightarrow\quad
\frac{1}{R_{\mathrm{eq}}}=\sum_k\frac{1}{R_k},
\qquad
R_{\mathrm{eq}}=\frac{R_1R_2}{R_1+R_2}\ \text{(two resistors)}.
\]
Series adds resistances; parallel adds conductances. That single distinction---same
current vs.\ same voltage---is the whole idea.

\section{Series and Parallel Capacitors}
The same two laws fix how capacitors combine, but the result comes out
\emph{backwards} from resistors---and it is worth seeing why rather than
memorizing it.

\textbf{Capacitors in series} carry the same current, so each accumulates the same
charge \(q\). Since \(v=q/C\) for each one, KVL sums the voltages:
\[
v=\frac{q}{C_1}+\frac{q}{C_2}+\cdots=q\sum_k\frac{1}{C_k}
\quad\Longrightarrow\quad
\frac{1}{C_{\mathrm{eq}}}=\sum_k\frac{1}{C_k},
\qquad
C_{\mathrm{eq}}=\frac{C_1C_2}{C_1+C_2}\ \text{(two capacitors)}.
\]
\textbf{Capacitors in parallel} share the same voltage, so their stored charges add
directly, \(q=C_1v+C_2v+\cdots\):
\[
C_{\mathrm{eq}}=\sum_k C_k .
\]

\begin{physicalbox}
\textbf{Why capacitors look ``backwards.''} Capacitance is charge stored per volt,
which makes it an ``ease'' quantity like conductance, not an ``opposition'' quantity
like resistance. So capacitance adds where conductance adds (parallel) and combines
reciprocally where resistance adds (series). The consistent way to see it is through
reactance: \(X_C=1/(\omega C)\) is the opposition, and \emph{reactances} in series
add exactly like resistors---which forces \(1/C\) to add. Two equal capacitors in
series give half the capacitance and twice the reactance, just as two equal
resistors in series give twice the resistance.
\end{physicalbox}

\section{Series and Parallel Inductors}
Inductors behave like resistors, because inductance is an opposition-like quantity
(\(X_L=\omega L\) grows with \(L\)). \textbf{In series} the same current flows
through each, so KVL adds their voltages \(L_k\,\dd i/\dd t\):
\[
L_{\mathrm{eq}}=\sum_k L_k .
\]
\textbf{In parallel} each sees the same voltage, so KCL adds their currents:
\[
\frac{1}{L_{\mathrm{eq}}}=\sum_k\frac{1}{L_k},
\qquad
L_{\mathrm{eq}}=\frac{L_1L_2}{L_1+L_2}\ \text{(two inductors)}.
\]
These assume no magnetic coupling between the coils; mutual inductance adds a cross
term.

\begin{center}
\begin{tabularx}{\textwidth}{@{}L{0.17\textwidth}Y Y@{}}
\toprule
Component & Series (same current) & Parallel (same voltage) \\
\midrule
Resistors \(R\) & \(R_{\mathrm{eq}}=\sum_k R_k\) & \(1/R_{\mathrm{eq}}=\sum_k 1/R_k\) \\
Inductors \(L\) & \(L_{\mathrm{eq}}=\sum_k L_k\) & \(1/L_{\mathrm{eq}}=\sum_k 1/L_k\) \\
Capacitors \(C\) & \(1/C_{\mathrm{eq}}=\sum_k 1/C_k\) & \(C_{\mathrm{eq}}=\sum_k C_k\) \\
\bottomrule
\end{tabularx}
\end{center}

\begin{workedbox}[: combining reactive parts]
Two \(\SI{220}{\pico\farad}\) capacitors in series give
\[
C_{\mathrm{eq}}=\frac{(220)(220)}{220+220}=\SI{110}{\pico\farad},
\]
half the value---the standard trick for making a smaller capacitance out of two
larger ones. In parallel the same pair gives \(\SI{440}{\pico\farad}\). Two
\(\SI{10}{\micro\henry}\) inductors do the opposite: \(\SI{20}{\micro\henry}\) in
series, \(\SI{5}{\micro\henry}\) in parallel.
\end{workedbox}

\begin{exambox}
The exam expects both directions from memory, and the reliable check is
reactance: whichever connection increases the total \emph{opposition} at a given
frequency is the one that behaves like series resistors. Series inductors add;
series capacitors combine reciprocally.
\end{exambox}

\section{Voltage and Current Dividers}
A voltage divider is just two series resistors read across one of them
(\cref{fig:divider}). The same current \(I=V_i/(R_1+R_2)\) flows through both, so
the output taken across \(R_2\) is
\[
V_o=IR_2=V_i\frac{R_2}{R_1+R_2}.
\]

\begin{figure}[htbp]
\centering
\begin{circuitikz}[scale=1]
\draw (0,0) to[V,l=\(V_i\)] (0,3)
      to[R,l=\(R_1\)] (2.5,3)
      to[R,l=\(R_2\),v=\(V_o\)] (2.5,0) -- (0,0);
\draw (2.5,3) to[short,-o] (3.5,3) node[right]{$+$};
\draw (2.5,0) to[short,-o] (3.5,0) node[right]{$-$};
\end{circuitikz}
\caption{Voltage divider: the output across \(R_2\) is
\(V_o=V_i R_2/(R_1+R_2)\).}
\label{fig:divider}
\end{figure}

Dually, current entering two parallel resistors splits \emph{inversely} with
resistance,
\[
I_1=I\frac{R_2}{R_1+R_2},\qquad I_2=I\frac{R_1}{R_1+R_2},
\]
so more current takes the lower-resistance branch.

\section{Thevenin and Norton Equivalents}
Any linear one-port of sources and resistors looks, from its terminals, like a
single Thevenin source \(V_{Th}\) in series with \(R_{Th}\), or equivalently a
Norton source \(I_N\) in parallel with \(R_N\) (\cref{fig:thevenin}):
\[
R_N=R_{Th},\qquad I_N=\frac{V_{Th}}{R_{Th}}.
\]
This matters in radio because a receiver input, an antenna, an amplifier stage, or
a test instrument is usually characterized by exactly this terminal
model---one source and one impedance---so you can predict how it behaves when a
different load or stage is connected without re-analyzing the whole circuit.

\begin{figure}[htbp]
\centering
\begin{circuitikz}
\draw (0,0) to[V,l=\(V_{Th}\)] (0,2)
      to[R,l=\(R_{Th}\)] (3,2)
      to[short,-o] (4,2);
\draw (0,0) -- (3,0) to[short,-o] (4,0);
\node at (4.8,1) {\(\Longleftrightarrow\)};
\draw (6,0) -- (6,2);
\draw (6,0) to[I,l_=\(I_N\)] (6,2);
\draw (6,2) to[R,l=\(R_N\)] (9,2) -- (9,0) -- (6,0);
\draw (9,2) to[short,-o] (10,2);
\draw (9,0) to[short,-o] (10,0);
\end{circuitikz}
\caption{Equivalent Thevenin and Norton one-port descriptions.}
\label{fig:thevenin}
\end{figure}

\section{Power Transfer and Matching}
For a purely resistive Thevenin source, maximum power reaches the load when
\[
R_L=R_{Th},
\]
and for complex impedances (developed in \cref{ch:ac}) the maximum-real-power
condition is the conjugate match \(Z_L=Z_S^{*}\). This is why RF systems are built
around a common reference
impedance: matching the source, line, load, and instruments to (typically)
\(\SI{50}{\ohm}\) delivers the most power to the antenna and, as \cref{ch:lines}
shows, minimizes reflections on the feed line.

\begin{exambox}
\(\SI{50}{\ohm}\) is a design convention, not a law of nature. Matching means the
system has been arranged so sources, lines, loads, and measuring instruments
interact predictably---maximum power transfer and minimum reflection---not that
every impedance is \(\SI{50}{\ohm}\).
\end{exambox}

\section{Worked Example}
A \(\SI{12}{\volt}\) source feeds \(R_1=\SI{2.2}{\kilo\ohm}\) on top and
\(R_2=\SI{1.0}{\kilo\ohm}\) to ground (the topology of \cref{fig:divider}). The
unloaded output and the Thevenin resistance seen by a load are
\[
V_o=12\frac{1.0}{2.2+1.0}=\SI{3.75}{\volt},
\qquad
R_{Th}=R_1\parallel R_2=\frac{(2.2)(1.0)}{2.2+1.0}\,\si{\kilo\ohm}
\approx\SI{0.688}{\kilo\ohm}.
\]
So to any load, this divider looks like a \(\SI{3.75}{\volt}\) source behind
\(\SI{688}{\ohm}\)---and a load comparable to \(\SI{688}{\ohm}\) will noticeably
pull the output down, which is why divider outputs are buffered before driving a
real load.
